% Generated draft: requires mathematical and source review.
\documentclass[11pt]{article}
\usepackage[margin=25mm]{geometry}
\usepackage{fontspec}
\setmainfont{Latin Modern Roman}
\usepackage{amsmath,amssymb,mathtools,unicode-math}
\setmathfont{Latin Modern Math}
\usepackage{enumitem,needspace,xurl,hyperref}
\setlength{\parindent}{0pt}
\setlength{\parskip}{5pt}
\setlength{\emergencystretch}{3em}
\newcommand{\N}{\mathbb{N}}
\newcommand{\Z}{\mathbb{Z}}
\newcommand{\Q}{\mathbb{Q}}
\newcommand{\R}{\mathbb{R}}
\newcommand{\C}{\mathbb{C}}
\newcommand{\F}{\mathbb{F}}
\newcommand{\A}{\mathbb{A}}
\newcommand{\PP}{\mathbb{P}}
\newcommand{\E}{\mathbb{E}}
\newcommand{\eps}{\varepsilon}
\newcommand{\dd}{\,\mathrm{d}}
\begin{document}
\section*{UnsolvedMath ID: 30003832}
\textbf{Source problem code:} catalogue entry
\subsection{Definitions and mathematical statement}
For each positive integer \(n\), let
\[
\mathcal M_n=\left\{X^n+\sum_{i=0}^{n-1}\varepsilon_iX^i:\varepsilon_i\in\{-1,1\}\right\}\subset\Z[X],
\qquad |\mathcal M_n|=2^n.
\]
For a monic polynomial \(f\) with complex roots \(\alpha_1,\ldots,\alpha_n\), counted with multiplicity, define
\[
\operatorname{Disc}(f)=\prod_{1\le i<j\le n}(\alpha_i-\alpha_j)^2\in\Z.
\]
There are two potentially relevant counting functions:
\[
S_n=\#\{f\in\mathcal M_n:\operatorname{Disc}(f)=a^2\text{ for some }a\in\Z\},
\]
\[
S_n^{\ne0}=\#\{f\in\mathcal M_n:\operatorname{Disc}(f)=a^2\text{ for some }a\in\Z\setminus\{0\}\}.
\]
The first includes discriminant zero and hence polynomials with repeated roots; the second excludes them. The supplied record does not specify whether zero is intended to count as a square or whether a separability or irreducibility restriction was intended. No such restriction is imposed here without qualification.

Under the interpretation that every integer square, including zero, is counted, the cleaned statement asks the following two questions:
\[
\lim_{n\to\infty}\frac{S_n}{2^n}=0\;?
\]
\[
\exists C>0\ \exists N\in\N\ \forall n\ge N,\qquad S_n\le C\;?
\]
The latter is the assertion \(S_n=O(1)\), with the constant independent of \(n\). If only nonzero square discriminants are intended, replace \(S_n\) by \(S_n^{\ne0}\) in both questions. The original statement instead asks whether the unnormalized number tends to zero. That would be the different assertion \(\lim_{n\to\infty}S_n=0\), with the same convention issue. The discrepancy between the original and cleaned statements remains unresolved; the displayed proportion question follows the cleaned statement provisionally. The original also refers to improving a relevant upper bound, but supplies neither that bound nor its hypotheses.

Separately, let \(A_5\) be the alternating group on five letters and let \(G_n=A_5^n\) be its direct product with \(n\) factors, provisionally for every positive integer \(n\). For a finite group \(H\), let \([H,H]\) be its commutator subgroup and let \(d(H)\) be the least cardinality of a generating set of \(H\), with \(d(\{1\})=0\). Since \(A_5\) is perfect, \(G_n/[G_n,G_n]\) is trivial. The supplied Boston--Markin prediction specializes to
\[
\forall n\ge1,\qquad m(G_n)=\max\{d(G_n/[G_n,G_n]),1\}=1.
\]
Here \(m\) denotes the minimal-ramification quantity mentioned in the source record. Its arithmetic definition is not supplied: the base field, admissible Galois realizations, and convention for counting ramified places are unspecified. Consequently this last display records the source's formal target, but does not yet specify a complete arithmetic conjecture.
\subsection{Short English statement}
Provisionally following the cleaned statement, ask whether the fraction of monic degree-\(n\) polynomials with coefficients in \(\{-1,1\}\) and square discriminant tends to zero, and whether their number is bounded independently of \(n\). Separately, ask whether the Boston--Markin minimal-ramification prediction gives \(m(A_5^n)=1\) for every positive integer \(n\). The counting normalization, treatment of zero discriminants, and definition of \(m\) require clarification.
\subsection{Sources}
\begin{itemize}
\item Supplied local UnsolvedMath registry record, permanent ID \(30003832\), exact source problem code \texttt{catalogue entry}, title: \emph{Square Discriminants and Minimal Ramification of Alternating Group Powers}. The record belongs to the Oberwolfach Reports Open Problems collection and supplies proposed year 2018; no individual proposer is named.
\item Supplied report citation: \emph{Field Arithmetic (2019)}, Oberwolfach Reports, pp.\ 1517--1577, DOI \texttt{10.4171/owr/2018/25}. \url{https://doi.org/10.4171/owr/2018/25}. The record describes these reports as documenting mathematical workshops and their open-problem discussions. The supplied 2019 label and the 2018 DOI component are retained without reconciliation.
\item Supplied literature source labeled ``Oberwolfach Report 2019'': \url{https://ems.press/content/serial-article-files/46747}.
\item Supplied literature source attributed to Bary-Soroker--Ben-Porath--Matei: \url{https://doi.org/10.1112/blms.13049}. No article title or precise supporting result is supplied.
\item Literature assessment checked at \texttt{2026-08-21}: supplied dated curation, not independently verified. It labels both the square-discriminant asymptotic and the \(A_5\)-product minimal-ramification conjecture open, and states that later work provides related bounds and lower bounds without resolving either target. The background explicitly requests independent verification of time-sensitive status and cited claims. The registry also supplies status ``open'' and statement status \texttt{corrected\_verified}; these are source labels, not independent mathematical verification in this draft.
\end{itemize}
\end{document}
